\include{header}
\begin{document}
\include{Glosario}
\include{Acronimos}
\include{portada}
\thispagestyle{empty}
\restoregeometry
\pagestyle{fancyplain}

\lhead{}\chead{}\rhead{}\lfoot{}\cfoot{}\rfoot{}
\include{agradecimientos}
\include{resumen}

%E: Página par
%O: Página impar
%L: Parte izquierda	
%C: Parte central
%R: Parte derecha
%H: Cabecera	
%F: Pie de página

\fancyhf{}
%pie de página, páginas impares izquierda, páginas pares derecha
\fancyfoot[LO,RE]{\fancyplain{}{\small\textit{Nombre del estudiante}}}

%pie de página, páginas pares izquierda, páginas impares derecha
\fancyfoot[LE,RO]{\fancyplain{}{\small\textit{Instituto Tecnológico Superior de Huichapan}}}

%cabecera de página, páginas impares izquierda, páginas pares derecha
\fancyhead[LO,RE]{\small\textit{ }}

%cabecera de página, páginas pares izquierda, páginas impares derecha
\fancyhead[LE,RO]{\fancyplain{}{\small\textit{\nouppercase{\leftmark}}}}

%número de página al centro
\fancyfoot[C]{\thepage}

\renewcommand{\headrulewidth}{0.6pt}
\renewcommand{\footrulewidth}{0.6pt}

\setlength{\parindent}{0.6cm}

\setcounter{page}{1}
\pagenumbering{roman}

%\include{resumen}
\renewcommand\contentsname{Contenido}
\tableofcontents
\renewcommand\listfigurename{Lista de figuras}
\listoffigures
\renewcommand\listtablename{Lista de tablas}
\listoftables

\renewcommand*{\glossaryname}{Glosario}
\renewcommand*{\acronymname}{Acrónimos}
\deftranslation{Glossary}{Glosario}
\deftranslation{Acronyms}{Acrónimos}
\printglossaries

\renewcommand\lstlistingname{Código}
\renewcommand\lstlistlistingname{Código}
\renewcommand\appendixname{Anexo}

%\newpage
\include{chapter1}
\include{chapter2}
\include{chapter3}
\include{chapter4}
\include{chapter5}

%insert the bibliography
\renewcommand{\bibname}{Referencias}
\bibliographystyle{unsrt}
\bibliography{main}
	
%\include{anexo1}

\appendix
\chapter{Anexo A: título}
\section{Anexo A: sección 1}
\noindent
Información importante para el lector pero que no es relevante para el marco teórico \cite{edo_art_14}.

\noindent
Incluir una imagen.

\begin{figure}[!h]
  \centering
    \includegraphics[width=0.4\textwidth]{images/Capitulo2/ciclo_de_vida_android_1.png}
  \caption[Ciclo de vida de Android.]{Ciclo de vida de Android \cite{W_is_android}.}
  \label{ciclo_vida_android}
\end{figure}

\subsection{Anexo A: sección 1.1.1}
Incluir una tabla.

\begin{table}[htp]
	\caption{Operadores relacionales}
	\begin{center}
		\begin{tabular}{|| c | l | c | l ||}
			\hline
			\hline
			Operador & nombre & ejemplo & significado\\
			\hline
			$<$ & menor que & a$<$b & a es menor que b\\
			\hline
			$>$ & mayor que & a$>$b & a es mayor que b\\
			\hline
			== & igual a & a==b & a es igual a b\\
			\hline
			!= & no igual a & a!=b & a no es igual a b\\
			\hline
			$<$= & menor que o igual a & a$<$=5 & a es menor que o igual a b\\
			\hline
			$>$= & mayor que o igual a & a$>$=b & a es menor que o igual a b\\
			\hline
		\end{tabular}
	\end{center}
\label{table_2-6}
\end{table}

\begin{lstlisting}[language=ListCpp,
caption={Firmware para Arduino de la prueba de integración de módulos externos.},
label = {codigo1}]
#include<SPI.h>
#include<Expander32.h>
#include<SpiMotorsV4.h>
#include<MCP23S17.h>
#include<MCP3208.h>
#include<AD8804.h>

SpiMotorsV4 motor(11);
MCP in(0,9);
MCP out(0,32);
MCP3208 adc(15);
AD8804 dac(26);

void setup(){
  Slave.begin(40,41);
  motor.setQEI(MODE_X4_RESET_ON_MATCH,VELOCITY_PULSE_DIV_1,DIV_BY_1,4800);
  out.pinMode(5,LOW);
  out.pinMode(6,LOW);
  out.pinMode(7,LOW);
  out.pinMode(8,LOW);
  out.digitalWrite(5,LOW);
  out.digitalWrite(7,LOW);
  in.pinMode(13,HIGH);
  in.pinMode(16,HIGH);
  Serial.begin(9600);
}

boolean state=false;
float speed = 0;
float temp = 0;
void loop(){
  state=!state;
  digitalWrite(13,state);
  out.digitalWrite(8,state);
  speed = motor.getSpeed();
  temp = adc.analogRead(CH0)* 0.122100122;
  if (in.digitalRead(13)){
    dac.analogWrite(A1,230);  //B
    dac.analogWrite(A2,120);  //G
    dac.analogWrite(A3,150);  //R
    out.digitalWrite(6,HIGH);
    Serial.print("temp= ");
    Serial.print(temp);
    Serial.print(" C\t");
    Serial.print("speed= ");
    Serial.print(speed);
    Serial.println("\t");
  }else if (in.digitalRead(16)){
    out.digitalWrite(6,LOW);
    dac.analogWrite(A1,0);
    dac.analogWrite(A2,0);
    dac.analogWrite(A3,255);
  }else {
    out.digitalWrite(6,LOW);
    dac.analogWrite(A1,160);
    dac.analogWrite(A2,120);
    dac.analogWrite(A3,255);
  }
}
\end{lstlisting}
\end{document}